\chapter{The Operating Model \& Startup Projections}\label{ch:operating-model}

% Scope: unit-economics-to-P&L bridge (cohort revenue build), the three-driver EBITDA model,
%   scenario budgeting & sensitivity, runway. NEW chapter (closes the financial-modeling gap).
% Bridges unit economics (\cref{ch:customer-economics}) to the DCF (\cref{ch:valuation}):
%   builds the FCF projections the DCF takes as given.
% Cross-links: \cref{ch:customer-economics} (CLV/CAC), \cref{ch:valuation} (DCF),
%   \cref{ch:fundraising-and-dilution} (runway trigger), \cref{ch:financial-statements} (three-statement).
% Draft per house style: flowing prose + example/admonition set-offs; running-SaaS worked examples.

\section{The Unit-Economics-to-P\&L Bridge}\label{sec:unit-econ-to-pl}
% complexity: 3

How does a pile of unit-economics assumptions turn into a revenue forecast? The answer is the
cohort stack: each month's new customers enter a cohort that shrinks by attrition each
subsequent month, contributing to ARR for as long as they remain active. Add up the surviving
contributions of every cohort and you have the firm's current ARR. This is not a top-down
TAM-share exercise; it is a bottom-up accounting identity that makes every revenue projection
traceable to two operational levers --- new-customer inflow and retention --- and one pricing
lever, ARPU.

Let $N_s$ be the number of new customers acquired in month $s$, ARPU the annualised revenue per
customer, and $c$ the monthly churn rate. The ARR at the end of period $t$ is the sum over all
prior cohorts of those customers still active:
\begin{equation}\label{eq:cohort-revenue}
  \mathrm{ARR}_{t}
  \;=\;
  \mathrm{ARPU}
  \sum_{s=1}^{t}
  N_{s}\,(1-c)^{\,t-s}.
\end{equation}
When $N_s$ is constant at $N$ and the series has been running long enough to approach steady
state, the sum collapses to $N\,\mathrm{ARPU}/c$; the exact steady-state ceiling is reached only
asymptotically, so real firms are always somewhere on the growth curve below it.
\Cref{eq:cohort-revenue} is the revenue twin of the CLV formula: \cref{eq:clv} values a single
customer, while \cref{eq:cohort-revenue} stacks the same survival mechanics across every cohort
in the book. The MRR growth rate identity in \cref{eq:mrr-growth} is the derivative of
\cref{eq:cohort-revenue} with respect to time.

\begin{example}[Cohort stack to ARR and Year-1 FCF]
The firm acquires \num{200} new customers per month at ARPU \money{600}/yr, with monthly churn
\qty{0.65}{\percent}. After \num{38} months of ramp, \cref{eq:cohort-revenue} yields:
\[
  \mathrm{ARR}_{38}
  \;=\; \money{600}\times
        \frac{200\,\bigl(1-0.9935^{38}\bigr)}{0.0065}
  \;\approx\; \money{4050000},
\]
which rounds to the \money{4000000} ARR in the operating plan. At this base, with invested
capital \money{3000000} and ROIC \qty{20}{\percent} (\cref{eq:roic}), NOPAT equals
\(\money{3000000}\times 0.20=\money{600000}\). Because a SaaS firm carries near-zero capital
expenditure and minimal working-capital swings, free cash flow is approximately equal to NOPAT;
\emph{Year-1 FCF \money{600000}} feeds directly into the DCF in \cref{eq:dcf} as the anchor
of the explicit forecast period, whose present-value sum is approximately \money{3377825}.

The implication is operational: both the \money{4000000} ARR and the \money{600000} FCF are
downstream outputs of three observable rates --- inflow, retention, and price. A founder who
cannot name the current value of all three cannot close the loop from marketing spend to
firm value.
\end{example}

\begin{admonition}{Warning}
Blended retention statistics hide cohort heterogeneity. A company that acquires half its
customers through a low-retention channel will report acceptable aggregate churn while the
healthy cohorts subsidise the leaky ones. Model retention \emph{by acquisition cohort}; the
top-of-funnel mix is as important as the rate itself.
\end{admonition}

The cohort stack connects directly to the P\&L: ARR is the revenue line; the variable cost
structure from \cref{ch:financial-statements} --- gross margin \qty{70}{\percent}, contribution
margin \qty{42}{\percent} --- determines what fraction of each cohort's revenue flows to
EBITDA. The next section turns those rates into an explicit driver model.


\section{The Three-Driver Model}\label{sec:three-driver-model}
% complexity: 3

Once ARR is projected from the cohort stack, the P\&L follows from three rates: gross margin,
the fixed-cost base, and its composition across the three operating expense buckets. Gross
margin reflects the engineering of the product and cannot be changed quickly; the cost
composition reflects management choices about where to invest. The EBITDA line is nothing but
the arithmetic residual of these rates applied to a growing ARR base. The discipline of
expressing the model this way is that it forces the team to be explicit about when operating
leverage kicks in --- the point at which ARR growth outpaces the growth of the fixed-cost
base and EBITDA turns positive.

Formalise the relationship as:
\begin{equation}\label{eq:op-model}
  \mathrm{EBITDA}_{t}
  \;=\;
  \mathrm{ARR}_{t}\cdot g_{\mathrm{m}}
  \;-\;
  \mathrm{FixedCosts}_{t},
\end{equation}
where $g_{\mathrm{m}}$ is the gross margin rate and $\mathrm{FixedCosts}_{t}$ is the sum of
Sales \& Marketing (S\&M), Research \& Development (R\&D), and General \& Administrative (G\&A)
costs. In a SaaS business the gross margin is relatively stable (the cost of serving one more
customer at scale is near zero), so EBITDA expansion is primarily a function of how quickly
ARR growth outpaces fixed-cost growth --- the textbook definition of operating leverage.
\textcite{koller2025valuation} trace this dynamic through the NOPAT bridge; \textcite{ross2022}
give the corporate-finance treatment of operating leverage and its implications for equity risk.

\begin{example}[Three-driver progression: EBITDA margin and Rule of 40]
Starting from \money{4000000} ARR growing at \qty{30}{\percent} per year, with gross margin
\qty{70}{\percent}, the operating expense mix evolves as the firm scales:

\medskip
\noindent\resizebox{\linewidth}{!}{%
\begin{tabular}{lrrrr}
\toprule
 & Year~0 & Year~1 & Year~2 & Year~3 \\
\midrule
ARR              & \money{3077000} & \money{4000000} & \money{5200000} & \money{6760000} \\
Gross profit (70\%) & \money{2154000} & \money{2800000} & \money{3640000} & \money{4732000} \\
S\&M (\% of ARR) & 50\%       & 50\%       & 45\%       & 35\%       \\
R\&D (\% of ARR) & 25\%       & 25\%       & 25\%       & 20\%       \\
G\&A (\% of ARR) & 10\%       & 10\%       & 15\%\textsuperscript{*} & 8\%        \\
Total opex       & 85\%       & 85\%       & 85\%       & 63\%       \\
EBITDA margin    & \(-15\%\)  & \(-15\%\)  & \;\;0\%    & \;\,+7\%   \\
\bottomrule
\end{tabular}%
}

\medskip
\noindent\textsuperscript{*}G\&A rises temporarily in Year~2 as the firm invests in finance
and legal infrastructure ahead of a fundraise; it compresses to \qty{8}{\percent} in Year~3
as ARR scales past the G\&A step function.

\medskip
The Rule-of-40 score (\cref{eq:rule40}) is the sum of ARR growth rate and EBITDA margin: Year~1
scores \(30+(-15)=15\), Year~2 scores \(30+0=30\), and Year~3 scores \(30+7=37\) --- a
progression from deeply unprofitable to approaching the benchmark threshold. The implication:
the operating-leverage story is only credible if the G\&A and R\&D curves genuinely flatten as
ARR scales. Any model that holds opex percentages constant as ARR grows is not modelling
operating leverage; it is modelling a cost-plus business.
\end{example}

\begin{admonition}{Warning}
Operating leverage cuts both ways. The same fixed-cost base that produces EBITDA expansion at
\qty{30}{\percent} growth produces accelerating losses if growth slows to \qty{15}{\percent}.
Stress-test the model at \qty{50}{\percent} of planned growth before presenting it to a board:
at \qty{15}{\percent} ARR growth the Year-3 EBITDA margin in the example above stays near
\(-5\%\) and the Rule-of-40 score does not cross the threshold. A plan that requires full
execution to avoid a cash crisis is a plan without margin for error.
\end{admonition}

\Cref{eq:op-model} outputs the EBITDA that feeds into the NOPAT bridge (\cref{eq:nopat-bridge})
and ultimately into the free cash flow that \cref{eq:dcf} discounts. The next section asks how
sensitive that output is to the assumptions driving it.


\section{Scenario Budgeting \& Sensitivity}\label{sec:scenario-budget}
% complexity: 4

A single-scenario operating model is a forecast in disguise. The honest version of the same
work is a sensitivity surface that maps the model's output across a range of input assumptions,
so management knows which lever matters most and by how much. The two most dangerous drivers in
a SaaS unit-economics model are CAC and churn: CAC governs the cost of filling the top of the
cohort stack; churn governs the speed at which it empties. Both can move independently, and
their interaction is non-linear.

The partial derivative of EBITDA with respect to CAC (holding all else constant) is simply the
number of new customers acquired per period. At \num{200} new customers per month:
\begin{equation}\label{eq:sensitivity}
  \begin{aligned}
    \frac{\partial\,\mathrm{EBITDA}}{\partial\,\mathrm{CAC}} \;=\; -N
    \quad&\Longrightarrow\quad \Delta\mathrm{CAC}=\money{100} \\
    &\Rightarrow\; \Delta\mathrm{EBITDA} = -\money{100}\times 200 = -\money{20000}\text{ per month},
  \end{aligned}
\end{equation}
or \money{240000} per year. The minus sign is load-bearing: CAC is an outflow that hits EBITDA
before a single dollar of lifetime value is realised. Churn operates through the ARR line via
\cref{eq:cohort-revenue}: a churn rate rise of \qty{0.35}{\percent} per month (from
\qty{0.65}{\percent} to \qty{1.00}{\percent}) reduces the steady-state customer count by roughly
\qty{35}{\percent}, which at \money{4000000} ARR is a structural \money{1400000} reduction in the
revenue ceiling --- an order of magnitude larger than a single year of CAC overspend.

\begin{figure}[htbp]
  \centering
  \includegraphics[width=0.92\textwidth]{op-model-bridge}
  \caption{Operating-model waterfall: from unit-economics inputs (ARPU, gross margin, churn,
    new customers per month) through ARR and EBITDA to free cash flow. The bar heights show
    how each driver contributes or subtracts; the rightmost bar is Year-1 FCF \money{600000}.}
  \label{fig:op-model-bridge}
\end{figure}

\begin{example}[Two-way sensitivity: 3-year NPV across CAC and churn]
The table below holds ARR\textsubscript{0} at \money{4000000}, annual new-customer
acquisition at \num{2400}, R\&D\,+\,G\&A at \money{2000000}, gross margin at \qty{70}{\percent},
and the discount rate at \qty{11}{\percent}. Each cell is the three-year net present value of
NOPAT, with a \qty{30}{\percent} tax rate applied to positive EBITDA.

\medskip
\noindent\resizebox{\linewidth}{!}{%
\begin{tabular}{lrrr}
\toprule
\multicolumn{1}{c}{CAC}
  & \multicolumn{1}{c}{Churn 0.50\%/mo}
  & \multicolumn{1}{c}{Churn 0.75\%/mo}
  & \multicolumn{1}{c}{Churn 1.00\%/mo} \\
\midrule
\money{400} & \money{2410000}  & \money{2120000}  & \money{1840000} \\
\money{500} & \money{2000000}  & \money{1710000}  & \money{1430000} \\
\money{600} & \money{1590000}  & \money{1300000}  & \money{1020000} \\
\money{700} & \money{1170000}  & \money{860000}   & \money{560000}  \\
\money{800} & \money{700000}   & \money{380000}   & \money{90000}   \\
\bottomrule
\end{tabular}%
}

\medskip
The base case (CAC \money{600}, churn \qty{0.65}{\percent}) produces a three-year NPV near
\money{1400000}, consistent with the \money{3377825} explicit-period PV in the full five-year
DCF (\cref{eq:scenario-ev}). The bear case --- CAC \money{800} and monthly churn
\qty{1.00}{\percent} --- reduces NPV to just \money{90000}, a \qty{94}{\percent} collapse
from base, illustrating that the two drivers compound: higher CAC raises the cost side at
exactly the moment higher churn is compressing the revenue base. Note that no cell in the table
is negative; even the bear case generates positive cumulative FCF over three years. The risk
at \money{800}/\qty{1.00}{\percent} is therefore not insolvency on this time horizon but
insufficient FCF to justify the valuation multiple required for a fundraise.
\end{example}

\begin{admonition}{Warning}
A model that is sensitive to only one driver is miscalibrated. If toggling CAC by
\qty{30}{\percent} changes NPV by \qty{50}{\percent} while toggling churn by the same
fraction changes it by only \qty{5}{\percent}, one of the two inputs is not being modelled with
enough granularity. Sensitivity analysis is not just a risk-communication tool; it is a
diagnostic that reveals whether the model's structure matches the actual drivers of the
business.
\end{admonition}

The sensitivity surface feeds the scenario logic. Three internally consistent scenarios ---
bear, base, and bull --- should each carry their own EBITDA trajectory and their own implied
runway, because the question the board asks after seeing the sensitivity table is always the
same: how long before we need to raise?


\section{Runway \& the Financing Decision}\label{sec:runway}
% complexity: 2

Runway is the bridge between operating model and financing decision. It converts the net cash
consumption rate into a countdown clock and sets the trigger date for fundraising. The
mechanics are simple; the discipline is in running two versions simultaneously --- a growth
runway that assumes planned new-customer acquisition continues, and a conservative runway that
assumes zero growth from the current base.

Let $B$ be the current cash balance, $E$ the monthly operating expenditure (gross burn), and
$R_{\mathrm{m}}$ the monthly cash actually collected from customers (MRR collected, not
deferred):
\begin{equation}\label{eq:runway}
  \text{Runway (months)}
  \;=\;
  \frac{B}{\,E - R_{\mathrm{m}}\,}
  \;=\;
  \frac{B}{\text{net monthly burn}}.
\end{equation}
The denominator is net burn: what the bank account loses each month after accounting for cash
inflows. The distinction between gross and net burn matters because a fast-growing firm with
substantial MRR has a very different survival clock than its income statement might suggest.
Gross burn is the baseline cost floor; net burn is the existential metric.
\textcite{brealey2020} develop the cash-flow statement mechanics that underpin the relationship
between net income, non-cash charges, and actual cash movement --- the same mechanics that
explain why MRR collected can differ from recognised revenue in the month it is booked.

\begin{example}[Runway and fundraising trigger]
The firm holds \money{2500000} in cash. Monthly operating expenditure (salaries, cloud
infrastructure, S\&M outflows) is \money{650000}; MRR currently collected is \money{340000},
consistent with the \money{4000000} ARR base and the fact that annual subscribers pay monthly.
Applying \cref{eq:runway}:
\[
  \text{Net burn}
  \;=\; \money{650000} - \money{340000}
  \;=\; \money{310000}\text{ per month},
\]
\[
  \text{Runway}
  \;=\; \frac{\money{2500000}}{\money{310000}}
  \;\approx\; \text{8 months}.
\]
The \emph{fundraising trigger} is the date that leaves enough runway to close a round. A
competitive Series A process takes four to six months from first outreach to wire; adding two
months of buffer sets the trigger at six months of runway remaining. Eight months of runway
minus six months of required lead time means the firm must begin fundraising \emph{immediately}
--- not at the six-month mark.

The implication for the operating model: the runway calculation should be re-run monthly
against actual cash and actual net burn. A burn rate that runs \money{50000} per month above
plan for three months has consumed a full round of buffer without a line appearing on the
income statement. The chapter on fundraising and dilution (\cref{ch:fundraising-and-dilution})
takes this trigger date and works through the valuation mechanics and term-sheet implications
of raising from this position.
\end{example}

\begin{admonition}{Warning}
Always run both the growth and the conservative runway, and use the conservative one to
set the trigger date. The growth runway is the upside scenario; the conservative runway ---
which holds MRR flat at today's level and ignores new sales --- is the survival clock. A firm
that runs its fundraising timing off the growth runway is betting its survival on the plan
executing exactly as modelled. The fundraising process will not wait for a recovery.
\end{admonition}

The operating model is now complete as a planning instrument: \cref{eq:cohort-revenue} builds
the ARR from observable unit-economics inputs; \cref{eq:op-model} converts ARR and cost
structure into EBITDA; \cref{eq:sensitivity} quantifies the impact of individual driver changes;
and \cref{eq:runway} converts the resulting cash profile into a survival timeline. The FCF
output of this chain --- \money{600000} in Year~1, growing at the rates in \cref{eq:op-model}
--- is the input that the DCF in \cref{eq:dcf} discounts to an enterprise value.
\textcite{koller2025valuation} close the loop between operating-model outputs and DCF
enterprise value in the chapters on forecasting and value drivers.
